\subsection{李群上点分布的卷积}

定义 1. 设 G 为李群，g 与 \(g'\) 为 G 的两点，且设 \(t \in \mathrm{T}_{g}^{(\infty)}(\mathrm{G})\)、\(t' \in \mathrm{T}_{g'}^{(\infty)}(\mathrm{G})\) 为 G 上在 g 与 \(g'\) 处的两个点分布 (Differentiable and Analytic Manifolds, R, 13.2.1)。t 与 \(t'\) 的卷积，记作 \(t * t'\)，是 \(t \otimes t'\) 在 G × G 到 G 的映射 \((h, h') \mapsto hh'\) 下的像 (Differentiable and Analytic Manifolds, R, 13.2.3)。

命题 1. (i) 若 \(t \in \mathrm{T}_g^{(s)}(\mathrm{G})\) 且 \(t' \in \mathrm{T}_{g'}^{(s')}(\mathrm{G})\)，则 \(t * t' \in \mathrm{T}_{gg'}^{(s + s')}(\mathrm{G})\)。

\begin{enumerate}
\def\labelenumi{(\roman{enumi})}
\setcounter{enumi}{1}
\tightlist
\item
  若 \(t\) 或 \(t'\) 没有常数项，则 \(t * t'\) 没有常数项。
\end{enumerate}

\[
\varepsilon_ {g} * \varepsilon_ {g ^ {\prime}} = \varepsilon_ {g g ^ {\prime}}.
\]

\begin{enumerate}
\def\labelenumi{(\roman{enumi})}
\setcounter{enumi}{3}
\tightlist
\item
  设 \(t \in \mathrm{T}_g^{(s)}(\mathrm{G})\)，\(t' \in \mathrm{T}_{g'}^{(s')}(\mathrm{G})\)，并设 \(f\) 是取值于一个分离多范空间的 \(\mathbf{C}^{s + s'}\) 类函数，定义在 \(gg'\) 的一个开邻域上。则
\end{enumerate}

\[
\begin{array}{c} \langle t * t ^ {\prime}, f \rangle = \langle t ^ {\prime}, h ^ {\prime} \mapsto \langle t, h \mapsto f (h h ^ {\prime}) \rangle \rangle \\ = \langle t, h \mapsto \langle t ^ {\prime}, h ^ {\prime} \mapsto f (h h ^ {\prime}) \rangle \rangle . \end{array}
\]

这由 Differentiable and Analytic Manifolds, R, 13.4.1、13.2.3 与 13.4.4 得出。

设 K = R 或 C，且 G 是有限维的。于是 G 是局部紧的。若 t、\(t'\) 是点测度，则 \(t \star t'\) 的定义与 Integration, Chapter VIII, § 1 中的定义一致。我们以后将看到，测度的卷积与点分布的卷积是不必为点分布的分布的卷积的两个特殊情形。

设 \(\mathcal{T}^{(\infty)}(\mathrm{G})\) 为诸 \(\mathrm{T}_g^{(\infty)}(\mathrm{G})\)（\(g\in \mathrm{G}\)）的直和 (cf.~Differentiable and Analytic Manifolds, R, 13.6.1)。我们把 \(\mathcal{T}^{(\infty)}(\mathrm{G})\) 的卷积定义为 \(\mathcal{T}^{(\infty)}(\mathrm{G})\times \mathcal{T}^{(\infty)}(\mathrm{G})\) 到 \(\mathcal{T}^{(\infty)}(\mathrm{G})\) 的、是定义 1 的卷积之扩充的双线性映射。我们也把它记作 \(*\)。这样，\(\mathcal{T}^{(\infty)}(\mathrm{G})\) 便具有一个由诸 \(\mathcal{T}^{(s)}(\mathrm{G})\) 滤过的代数结构。子代数 \(\mathcal{T}^{(0)}(\mathrm{G}) = \bigoplus_{g\in \Omega}\mathrm{T}_g^{(0)}(\mathrm{G})\) 与 K 上群 G 的代数 \(\mathrm{K}^{\mathrm{(G)}}\) 认同。

命题 2. 代数 \(\mathcal{T}^{(\infty)}(\mathrm{G})\) 是结合的。它交换当且仅当 \(\mathrm{G}\) 交换。

设 \(t \in \mathcal{T}^{(\infty)}(G)\)，\(t' \in \mathcal{T}^{(\infty)}(G)\)，\(t'' \in \mathcal{T}^{(\infty)}(G)\)。于是 \(t * (t' * t'')\) 是 \(t \otimes t' \otimes t''\) 在映射 \((g, g', g'') \mapsto g(g'g'')\)（从 \(G \times G \times G\) 到 G）下的像，而 \((t * t') * t''\) 是 \(t \otimes t' \otimes t''\) 在映射 \((g, g', g'') \mapsto (gg')g''\)（从 \(G \times G \times G\) 到 G）下的像。因此 \((t * t') * t'' = t * (t' * t'')\)。类似地可以看出，若 G 交换，则 \(t * t' = t' * t\)。若卷积交换，则由命题 1 (iii)，G 交换。

命题 3. 若 \(t \in \mathcal{T}^{(\infty)}(G)\) 且 \(g \in G\)，则 \(\gamma(g)_*t = \varepsilon_g * t\)，\(\delta(g)_*t = t * \varepsilon_{g^{-1}}\)，(Int \(g)_*t = \varepsilon_g * t * \varepsilon_{g^{-1}}\)。特别地，\(\varepsilon_e\) 是 \(\mathcal{T}^{(\infty)}(G)\) 的单位元。

考虑图

\[
\mathrm{G} \xrightarrow {\Phi} \mathrm{G} \times \mathrm{G} \xrightarrow {\Psi} \mathrm{G}
\]

其中 \(\phi\) 是映射 \(h\mapsto (g,h)\)，\(\psi\) 是映射 \((h^{\prime},h)\mapsto h^{\prime}h\)。于是 \(\gamma (g) = \psi \circ \phi\)，因而 \(\gamma (g)_*t = \psi_*(\phi_*(t))\)。但 \(\phi_{*}(t) = \varepsilon_{g}\otimes t\)，因而 \(\psi_{*}(\phi_{*}(t)) = \varepsilon_{g}*t\)。对 \(\delta (g)_*t\) 的论证是类似的。最后，

\[
\operatorname{Int} g = \gamma (g) \circ \delta (g)
\]

因而 \((\operatorname{Int} g)_* = \gamma(g)_* \circ \delta(g)_*\)。

由此可以看出，对 \(t \in \mathrm{T}(\mathrm{G})\)，\(\varepsilon_g * t\) 与 \(t * \varepsilon_g\) 分别等于 \(gt\) 与 \(tg\)（在群 \(\mathrm{T}(\mathrm{G})\) 中算得）(§ 2, 第 2 节)。但应注意，对 \(t\)、\(t'\)（\(\mathrm{T}(\mathrm{G})\) 中的元素），§ 2 意义下的乘积 \(tt'\) 一般不同于 \(t * t'\)。

定义 2. 设 \(\mathbf{G}\) 为李群。\(\mathcal{T}^{(\infty)}(\mathbf{G})\) 中由支集含于 \(e\) 的分布组成的子代数记作 \(\mathbf{U}(\mathbf{G})\)。

这个代数由诸子空间

\[
\mathrm{U} _ {s} (\mathrm{G}) = \mathrm{U} (\mathrm{G}) \cap \mathcal {T} ^ {(s)} (\mathrm{G}) = \mathrm{T} _ {e} ^ {(s)} (\mathrm{G}).
\]

我们记 \(\mathrm{U}^{+}(\mathrm{G}) = \mathrm{T}_{e}^{(\infty)+}(\mathrm{G}),\mathrm{U}_{s}^{+}(\mathrm{G}) = \mathrm{U}^{+}(\mathrm{G})\cap \mathrm{U}_{s}(\mathrm{G})\) (cf.~Differentiable and Analytic Manifolds, R, 13.2.1)。回忆，\(\mathrm{U}_0(\mathrm{G})\) 与 K 认同，而 \(\mathrm{U}_1^+\) (G) 与切空间 \(\mathrm{T_e(G)}\) 认同。在 \(\mathrm{U}(G)\) 中，\(\mathrm{U}^{+}(\mathrm{G})\) 是与 \(\mathrm{U}_0(\mathrm{G})\) 互余的双边理想。

例。设 E 为一个被看作李群的完备可赋范空间。于是向量空间 U(E) 典范地与向量空间 TS(E) 认同 (Differentiable and Analytic Manifolds, R, 13.2.4)。设 \(m: \mathrm{E} \times \mathrm{E} \to \mathrm{E}\) 为 E 上的加法。于是

\[
m _ {*} \colon \mathsf {T S} (\mathrm{E} \times \mathrm{E}) \rightarrow \mathsf {T S} (\mathrm{E})
\]

等于 \(\mathsf{TS}(m)\) (Differentiable and Analytic Manifolds, R, 13.2.4)。于是，对 \(t\)、\(t'\)（\(\mathrm{U(E)} = \mathrm{TS(E)}\) 中的元素），像 \(t * t'\)（即对称张量积 \(t \otimes t'\) 在 \(m_*\) 下的像）是 \(\mathsf{TS}(m)(t \otimes t')\)。由 Algebra, Chapter IV, § 5, no. 6, Proposition 7，这个像正是乘积 \(tt'\)（在代数 \(\mathsf{TS(E)}\) 中的乘积）。这样，代数 \(\mathrm{U(E)}\) 便与代数 \(\mathsf{TS(E)}\) 认同。

命题 4. 考虑双线性映射 \((u,v)\mapsto u*v\)（相应地 \((u,v)\mapsto v*u\)），它从 \(\mathrm{U}(\mathrm{G})\otimes \mathbf{K}^{(\mathrm{G})}\) 到 \(\mathcal{T}^{(\infty)}(\mathrm{G})\) 中。相应的线性映射（从 \(\mathrm{U}(\mathrm{G})\otimes \mathbf{K}^{(\mathrm{G})}\) 到 \(\mathcal{T}^{(\infty)}(\mathrm{G})\)）是向量空间同构。

\(K^{(G)}\) 是诸 \(K\varepsilon_{x}\)（\(x\in G\)）的直和。另一方面，由命题 3，映射 \(u\mapsto u*\varepsilon_{g}\)（相应地 \(u\mapsto\varepsilon_{g}*u\)）是向量空间 \(U(G)=\mathcal{T}_{e}^{(\infty)}(G)\) 到向量空间 \(\mathcal{T}_{g}^{(\infty)}(G)\) 上的同构。最后，\(\mathcal{T}^{(\infty)}(G)\) 是诸 \(T_{g}^{(\infty)}(G)\)（\(g\in G\)）的直和。

设 X 为 \(C^{r}\) 类流形（\(r \geqslant \infty\)），且 \(x \in X\)。我们已经在向量空间 \(\mathcal{T}_{x}^{(\infty)}(\mathbf{X})\) 上定义了一个典范滤过，并定义了一个典范同构 \(i_{X,x}\)（从相伴分次向量空间到分次向量空间 \(\mathsf{TS}(\mathrm{T}_{x}(\mathbf{X}))\) 上）(Differentiable and Analytic Manifolds, R, 13.3.1)。特别地，设 \(\mathrm{T}_{e}(\mathrm{G}) = L\)；于是 \(i_{G,e}\) 是分次向量空间 gr \(\mathrm{U}(\mathrm{G})\) 到分次向量空间 TS(L) 上的同构。但 \(\mathrm{U}(\mathrm{G})\) 是一个滤过代数，由此我们在 gr \(\mathrm{U}(\mathrm{G})\) 上得到一个分次代数结构。

命题 5. 同构 \(i_{G,e} \colon \operatorname{gr} U(G) \to TS(L)\) 是代数同构。

设 \(p\) 为 U(G) × U(G) 到 U(G × G) 的映射 \((t, t') \mapsto t \otimes t'\)。设 \(c\) 为 U(G) × U(G) 到 U(G) 的映射 \((t, t') \mapsto t * t'\)。设 m 为 G × G 到 G 的映射 \((g, g') \mapsto gg'\)。于是由定义 1

\begin{enumerate}
\def\labelenumi{(\arabic{enumi})}
\tightlist
\item
\end{enumerate}

\[
c = m _ {*} \circ p.
\]

考虑下图

\[
\begin{array}{c} \operatorname{gr} \mathrm{U(G)} \times \operatorname{gr} \mathrm{U(G)} \xrightarrow {\operatorname{gr} (p)} \operatorname{gr} \mathrm{U(G \times G)} \xrightarrow {\operatorname{gr} (m _ {*})} \operatorname{gr} \mathrm{U(G)} \\ i _ {\mathrm{G}, e} \times i _ {\mathrm{G}, e} \Biggl \downarrow \\ \mathsf {T S} (\mathrm{L}) \times \mathsf {T S} (\mathrm{L}) \xrightarrow {q} \mathsf {T S} (\mathrm{L} \times \mathrm{L}) \xrightarrow {\mathsf {T S} (\mathrm{T} (m))} \mathsf {T S} (\mathrm{L}) \end{array}
\]

其中 q 是由 TS(L) × TS(L) 到 TS(L × L) 的典范同构导出的映射。由 Differentiable and Analytic Manifolds, R, 13.4.6 与 13.3.5，图中两个方块都是交换的。因而由 (1)，图

\[
\begin{array}{c} \operatorname{gr} U (G) \times \operatorname{gr} U (G) \xrightarrow {\operatorname{gr} (c)} \operatorname{gr} U (G) \\ ^ {i _ {G, e} \times i _ {G, e}} \Biggl \downarrow \qquad \qquad \qquad \qquad \qquad \qquad \qquad \Biggl \downarrow^ {i _ {G, e}} \\ T S (L) \times T S (L) \xrightarrow {T S (T (m)) \circ q} T S (L) \end{array}
\]

是交换的。现在，\(\mathbf{T}(m): \mathbf{L} \times \mathbf{L} \to \mathbf{L}\) 把 \((x, y)\) 映为 \(x + y\) (§ 2, 第 1 节, 命题 2 (ii))。由 Algebra, Chapter IV, § 5, no. 6, Proposition 7，\(\mathbf{TS}(\mathbf{T}(m)) \circ q\) 因而是代数 \(\mathbf{TS}(\mathbf{L})\) 的乘法。

